\chapter{相关研究基础}

\section{医学图像检测与分割研究的总体脉络}
医学图像分析的发展大体经历了从手工特征、浅层机器学习到深度学习，再到预训练大模型和提示学习的几个阶段。早期方法主要依赖阈值分割、区域生长、活动轮廓模型、形态学算子以及人工设计纹理特征来完成病灶或组织结构提取。这类方法的优点是可解释性强、实现简单，但它们通常只能适用于成像条件较为稳定、目标形态变化较小的任务。一旦遇到病理染色变化、扫描噪声增强或者目标边界模糊等问题，传统方法往往很难保持稳定性能。

随着卷积神经网络的成功应用，医学图像分析逐渐进入深度学习阶段。以 U-Net 为代表的编码器--解码器结构推动了像素级分割任务的快速发展，Mask R-CNN 及其变体则推动了实例级分割的发展。对于细胞核检测与分割任务而言，深度学习方法显著提高了边界恢复能力和复杂背景下的识别能力。然而，这一阶段的方法几乎都建立在“存在足够强监督标注”的前提之上，也因此把问题转移到了数据构建成本上。医学数据的标注瓶颈并没有因为模型变强而消失，反而由于模型更依赖大规模训练而变得更加突出。

近几年，视觉语言预训练和基础模型的发展又把医学图像分析推入新的阶段。研究重点开始从“如何在固定数据集上逼近最优精度”转向“如何减少对强监督标注的依赖”“如何在开放类别和跨域场景中泛化”以及“如何利用文本、图像示例和外部知识协助推理”。本文所关注的弱监督、零样本、图像提示和数字病理分析问题，本质上都处在这一新阶段的演进链条中。

\section{弱监督与无监督医学图像分割}

\subsection{弱监督分割的基本思想}
弱监督分割的基本目标是利用比像素级标注更廉价的监督信息近似恢复目标区域。常见弱监督信号包括图像级标签、点标注、框标注、涂鸦标注以及部分伪标签。不同监督形式对应不同的信息上限：图像级标签最廉价，但只告诉模型某类目标是否存在；点和框能提供更明确的空间约束，但边界信息仍然缺失；涂鸦比点和框更接近真实轮廓，但人工成本也更高。如何在“监督强度”和“标注成本”之间取得平衡，是弱监督分割长期关注的问题。

在自然图像领域，类激活图、区域提案、条件随机场后处理和自训练策略构成了图像级弱监督分割的主流技术路线。即先通过分类网络获得目标的粗激活分布，再利用后处理或迭代优化方式把粗热图转化为较完整的掩码。在医学图像场景中，这一思路同样具有启发性，但直接照搬并不总是有效。医学对象尤其是细胞核目标具有尺度小、密度高、粘连严重等特点，使得自然图像中的弱监督前景恢复策略很难直接实现高质量实例边界提取。

\subsection{病理图像中的特殊困难}
病理图像弱监督分割面临的困难具有鲜明的结构性。首先，细胞核等目标往往尺寸小、数量多，相互之间距离很近，一旦伪标签稍微膨胀，就会把多个目标错误地连在一起。其次，病理图像中的颜色和纹理受染色和扫描条件影响显著，不同样本之间的成像风格变化容易导致弱监督定位结果不稳定。再次，病理图像中存在大量与目标外观接近的背景结构，这些结构往往会在弱监督阶段被错误激活，造成伪标签污染。

因此，病理图像中的弱监督分割方法通常不能只停留在“恢复前景区域”层面，而必须进一步思考如何形成实例级结构感知。这也是为什么细胞核分割工作中常常会显式引入边界分支、距离变换分支或中心点约束等结构先验。类似地，本文后续针对密集细胞核场景中的漏检与重叠问题，也引入了基于伪标签迭代的自训练/蒸馏机制，以在弱监督信号有限的条件下逐步提升结构化定位质量。

\subsection{伪标签与自训练}
伪标签和自训练在弱监督与半监督分割中占据重要位置。其基本思想是先利用现有模型生成高置信预测，再把这些预测当作新监督信号继续训练模型。成功的关键在于：伪标签既不能过于保守，否则覆盖率不足；也不能过于激进，否则错误会在迭代中不断放大。在医学图像中，这个平衡更难把握，因为目标形态复杂且域偏移明显。

已有研究证明，自训练在医学图像中并非简单的“多跑几轮”策略，而需要与结构先验、分支解耦或多任务约束结合。对于本文的研究主线而言，伪标签与自训练并不是一个孤立技术点，而是一种贯穿多篇工作的核心思想：从弱监督分割中的伪标签迭代优化，到零样本检测后的目标域自适配（自训练/蒸馏），再到图像提示学习中的少样本引导，都能看到“先构造一个可接受起点，再通过数据反馈逐步提升”的共通逻辑。

\section{细胞核检测与实例分割研究}

\subsection{检测、分割与联合建模}
细胞核分析任务通常分为三类：核检测、语义分割和实例分割。核检测强调目标位置，适合计数和密度估计；语义分割强调前景区域覆盖，适合估计总体核面积；实例分割则要求区分每个独立核目标，最适合后续的形态统计和组织微环境分析。三类任务并非彼此割裂，而是存在显著内在联系。一个好的实例分割模型往往隐含具有检测能力，一个好的检测器也通常能为实例学习提供较稳定候选区域。

在病理图像中，许多方法试图将检测与分割联合建模。例如，有的工作先进行核中心点检测，再以中心点为引导恢复边界；有的工作同时预测前景图、边界图和距离图，以实现实例拆分；还有的工作直接通过 proposal-based 结构完成实例掩码预测。不同方法的差异在于对实例结构的建模方式不同，但它们都说明一件事：对于高密度细胞核任务，仅靠单一前景预测往往不够，必须引入额外结构信息。

\subsection{跨数据集与跨域问题}
细胞核数据集之间的差异很大。不同组织来源、不同切片质量、不同染色方案和不同扫描仪都会导致目标外观变化。一个在单一公开 benchmark 上表现优异的模型，进入真实跨域场景后往往会遇到性能下降。因此，单纯追求单数据集分数并不足以说明方法的实用性。如何在有限目标域支持样本甚至无支持样本条件下增强泛化能力，已经成为细胞核分析研究的重要方向。

这一点与本文后续关于跨域评测（domain shift）、少样本支持提示（few-shot prompting）以及“冻结主干、仅通过提示适配”的直测设定密切相关。作者之所以反复强调源域训练、目标域少量支持样本提示与不进行目标域大规模微调，是因为细胞核分析任务的实际挑战并不仅仅是“在一个数据集上做得更高”，而是“面对新域时还能否工作”。从学位论文角度看，这种问题意识也使得全文研究比单纯追求基准分数更具完整性。

\section{视觉语言预训练模型及其医学迁移}

\subsection{视觉语言预训练的基本范式}
视觉语言预训练模型通过大规模图文对构建跨模态共享语义空间。CLIP 通过图像编码器和文本编码器的对比学习训练获得强分类迁移能力，BLIP 在此基础上进一步增强生成与理解能力，而 GLIP 则尝试把语言提示直接融入目标检测框架，使检测器具备开放词汇识别能力。这一系列模型的重要意义在于，它们使图像理解不再局限于固定类别表，而可以借助文本描述扩展到未见对象。

在自然图像领域，开放词汇检测的价值已经得到广泛验证。但在医学图像领域，这一能力并不能直接迁移。首先，医学图像的许多概念在大规模通用图文语料中本来就稀缺；其次，医学对象往往依赖局部结构模式而非简单语义标签；再次，医学任务常常需要非常稳定的定位能力，而非仅仅粗粒度识别。因此，视觉语言模型迁移到医学场景时，需要配合新的提示机制、目标域自适应机制和可能的结构先验。

\subsection{开放词汇与零样本能力的局限}
尽管视觉语言模型具备开放词汇能力，但在医学任务中，单靠文本类名通常不足以支撑高质量检测。例如，同样是“细胞核”，不同病理环境下的核在大小、形状、染色情况和周边背景方面可能完全不同。简单提示无法表达这种差异，甚至会导致模型偏向自然图像语义空间中的错误对象。因此，零样本医学检测如果要真正成立，必须在提示构造和目标域适配上做更细致设计。

本文第三章正是在这一局限之上提出自动属性提示构造。其出发点不是否认视觉语言模型的价值，而是承认“医学类别语义不足”的现实，然后试图以属性、描述和自训练的方式弥补这一不足。从相关工作回顾角度看，这使得本文不是简单使用现成大模型，而是在其基础上提出新的迁移方法。

\section{提示学习的发展与图像提示研究}

\subsection{文本提示学习}
提示学习最早主要体现为对文本模板的设计。研究者通过改变 prompt phrasing 来影响模型对类别的理解，例如“a photo of a dog”和“this is a dog”在 CLIP 类模型中可能带来不同效果。随着研究发展，更多工作开始引入可学习 prompt token，以代替人工设计模板。对医学图像而言，文本提示学习具有明显吸引力，因为它不需要大规模改动模型结构，却可以借助预训练语义进行目标域适配。

不过，医学场景中的文本提示存在两个根本局限。第一，语言本身未必能精确表达细粒度视觉差异。第二，许多医学任务对形态和组织结构极为敏感，而这些信息天然更容易通过视觉示例表达，而非通过词语表达。因此，仅依赖文本 prompt 往往不足以满足复杂医学场景需求。

\subsection{图像提示与多模态提示}
图像提示的核心思想是让支持图像本身参与推理，以补充纯文本提示不足。已有工作在通用目标检测、分割和 few-shot 学习中表明，少量支持样本可以显著改善模型对新类别的感知能力。但如何让图像提示与原有视觉语言预训练结构兼容，仍是一个开放问题。若处理不当，图像提示可能破坏原有图文对齐能力，导致模型难以稳定收敛。

在医学场景下，图像提示的潜力更加明显。一个病理专家往往只需要看几张代表性示例就能理解某类目标的大致外观，而图像提示恰好模拟了这种“示例式知识传递”过程。本文第五章的视觉文本化工作，就是对这一思路的系统展开：既不放弃文本提示提供的语义框架，又引入图像提示补充细粒度结构信息。

\section{数字病理分析与可解释医学人工智能}

\subsection{从识别到量化}
数字病理研究正在从“自动识别病灶”向“自动量化病理表型”演化。前者关注的是目标是否被正确找出，后者关注的是目标识别后能否形成对医生有意义的统计指标。例如，肿瘤残留比例、免疫细胞密度、肿瘤边界结构变化以及细胞空间共定位模式，都属于后者关心的内容。

这种演化意味着，前端检测和分割结果的意义不再只体现为 benchmark 分数，而是要进一步考虑结果是否足以支持后续统计分析与临床解释。因此，一个看似精度很高但结构不稳定的模型，未必真的适合数字病理分析；相反，一个在临床关键指标上更稳定、更可解释的模型，可能更有实际应用价值。

\subsection{临床结局建模}
在肿瘤研究中，治疗后病理切片提供了非常丰富的预后信息。除了传统的病理缓解程度评估，肿瘤周围免疫细胞的数量、密度和空间分布也逐渐被认为与治疗反应和生存结局有关。AI 模型若能稳定识别并量化这些指标，就能为更细致的风险分层和个体化治疗提供辅助。

本文第六章的肺鳞癌工作正是建立在这一思想之上。它使得前几章偏方法论的研究，有了一个自然的应用收束点，也让整篇学位论文的价值不局限于算法创新，而延伸到临床解释和数字病理转化。

\section{本章小结}
本章系统回顾了本文后续研究所依赖的相关基础，包括弱监督与无监督医学分割、细胞核检测与实例分割、视觉语言预训练模型、提示学习以及数字病理分析。从这些工作的发展脉络可以看到，本文四项研究并不是偶然拼接，而是分别对应“监督不足”“语义不足”“跨域适配”和“临床解释”四类连续问题。基于这一认识，后续章节将依次展开具体方法与实验分析。
